%% ============================================================
%%  HGST Preprint — Overleaf-ready
%%  Submitted to: arXiv:hep-lat / q-bio.MN
%%  Date: March 4, 2026
%% ============================================================
\documentclass[12pt,a4paper]{article}

%% --- Packages ---
\usepackage[T1]{fontenc}
\usepackage[utf8]{inputenc}
\usepackage{lmodern}
\usepackage{amsmath,amssymb,amsthm}
\usepackage{mathtools}
\usepackage{graphicx}
\usepackage{xcolor}
\usepackage{booktabs}
\usepackage{multirow}
\usepackage{array}
\usepackage{tabularx}
\usepackage{longtable}
\usepackage{hyperref}
\usepackage{cleveref}
\usepackage{geometry}
\usepackage{setspace}
\usepackage{caption}
\usepackage{subcaption}
\usepackage{tikz}
\usepackage{pgfplots}
\pgfplotsset{compat=1.18}
\usepackage{listings}
\usepackage{physics}
\usepackage{siunitx}
\usepackage{natbib}
\usepackage{microtype}
\usepackage{float}
\usepackage{enumitem}
\usepackage{mdframed}

\geometry{
  top=2.5cm, bottom=2.5cm,
  left=3.0cm, right=3.0cm
}
\onehalfspacing

%% --- Hyperref setup ---
\hypersetup{
  colorlinks=true,
  linkcolor=blue!70!black,
  citecolor=green!60!black,
  urlcolor=blue!60!black,
  pdftitle={Frustration Ordering in HGST: U(1) to SU(3)xSU(2)xU(1)},
  pdfauthor={HGST Development Team},
}

%% --- Custom commands ---
\newcommand{\R}{\mathcal{R}}           % MIXED fraction
\newcommand{\Rinf}{R_\infty}
\newcommand{\hgst}{\textsc{hgst}}
% \comm, \norm, \ket provided by physics package
\newcommand{\SU}[1]{\mathrm{SU}(#1)}
\newcommand{\UN}[1]{\mathrm{U}(#1)}
\newcommand{\SM}{\SU{3}\!\times\!\SU{2}\!\times\!\UN{1}}
\newcommand{\validated}{\textcolor{green!55!black}{\textbf{[VERIFIED]}}}
\newcommand{\supported}{\textcolor{blue!70!black}{\textbf{[SUPPORTED]}}}
\newcommand{\falsified}{\textcolor{red!70!black}{\textbf{[FALSIFIED]}}}
\newcommand{\open}{\textcolor{orange!70!black}{\textbf{[OPEN]}}}

%% --- Theorem environments ---
\theoremstyle{plain}
\newtheorem{theorem}{Theorem}[section]
\newtheorem{lemma}[theorem]{Lemma}
\newtheorem{proposition}[theorem]{Proposition}
\theoremstyle{definition}
\newtheorem{definition}[theorem]{Definition}
\theoremstyle{remark}
\newtheorem{remark}[theorem]{Remark}

%% --- Epistemic-status box ---
\newmdenv[
  linecolor=gray!60,
  backgroundcolor=gray!8,
  roundcorner=4pt,
  innertopmargin=6pt,
  innerbottommargin=6pt,
  skipabove=8pt,
  skipbelow=8pt
]{statusbox}

%% ============================================================
\begin{document}
%% ============================================================

%% --- Title block ---
\title{%
  \textbf{Frustration Universality in Hierarchical Graded Structure Theory:\\ 
  A Lattice Study of the E7 MIXED Class from U(1) to the $\mathrm{SU}(3){\times}\mathrm{SU}(2){\times}\mathrm{U}(1)$ Gauge Group}
}

\author{%
  HGST Development Team\\[4pt]
  \small\textit{Correspondence:} \href{mailto:hgst-team@example.org}{hgst-team@example.org}
}

\date{March 4, 2026 \quad
      {\normalsize(Preprint v1.0 --- Fifth Edition results)}}

\maketitle
\thispagestyle{empty}

%% ============================================================
\begin{abstract}
\noindent
We present a systematic lattice study of the \emph{MIXED frustration fraction} $\R$
in Hierarchical Graded Structure Theory (\hgst) across a progression of gauge groups:
$\UN{1}$ (E1), $\SU{2}$ (E7-SU2), $\SU{3}$ (E7-SU3), and the SM
product gauge group $\SM$ (E7-SM, classical scalar lattice).
The observable $\R$ counts the fraction of signed triads whose mediated sign is
inconsistent with the direct sign---a purely combinatorial measure of phase frustration.
For classical $\UN{1}$ gauge fields the Abelian commutator $[\,\cdot\,,\,\cdot\,]=0$ forces
$\R\to 0$, falsifying the naive HGST prediction.
Non-Abelian gauge groups systematically sustain $\R>0$: we find
$R_\infty^{\text{SU2}}=0.3669\pm0.0036$ and $R_\infty^{\text{SU3}}=0.3539\pm0.0195$ from
finite-size scaling (FSS) on $L=4,6,8,10$ lattices (autocorrelation-corrected errors).
The SU(3) value $R_\infty^{\text{SU3}}=0.3539\pm0.0195$ overlaps numerically
with the RegulonDB \textit{E.~coli} K-12 MIXED fraction
$\R_\text{RegulonDB}=0.349\pm0.018$ (Appendix~\ref{app:regulondb}); assessing
whether this numerical overlap reflects shared algebraic structure requires further
empirical testing (\S\ref{sec:discussion}).
The $\SU{3}/\SU{2}$ frustration ratio is $1.176\pm 0.015$; the commutator norm
ratio is $\mathcal{C}(\SU{3})/\mathcal{C}(\SU{2})=1.421$ (Appendix~\ref{app:commutator}),
confirming the ordering $\R[\SU{2}]<\R[\SU{3}]$ but indicating a non-linear mapping
between $\mathcal{C}(G)$ and $\R$.
For the SM product group we implement five new simulation modules (32/32 unit tests),
execute $\beta_3$-scans, $\kappa$-scans, FSS at $L=4,6,8$, and a quark-vs-lepton scan.
The extrapolated SM values $R_\infty^\text{quark}\approx 0.493$ and
$R_\infty^\text{lepton}\approx 0.480$ both exceed the RegulonDB measured range and
converge toward $\R=0.5$.
Quarks systematically exceed leptons by $\Delta\R\approx 0.02$--$0.03$;
EW-only simulations ($\beta_3=0$) confirm this gap is intrinsic to $\SU{2}\times\UN{1}$
and not an artefact of colour.
These results establish a \emph{frustration ordering} by Lie algebra
rank: $\UN{1}<\SU{2}<\SU{3}$, with $R_\infty^{\SU{3}}\approx\R_\text{RegulonDB}$ (Appendix~\ref{app:regulondb}).
\end{abstract}

\noindent\textbf{Keywords:} lattice gauge theory, non-Abelian gauge groups, frustration
universality, MIXED class, SM gauge group, finite-size scaling,
signed network frustration, combinatorial frustration, gene regulatory networks, HGST

\vspace{6pt}
\noindent\textbf{MSC/PACS:} 11.15.Ha, 05.50.+q, 87.18.Vf, 02.20.Qs

\newpage
\tableofcontents
\newpage

%% ============================================================
\section{Introduction}
\label{sec:intro}
%% ============================================================

Understanding how local algebraic structure generates global emergent complexity is a
central challenge spanning lattice gauge theory~\citep{Wilson1974}, network
science~\citep{Barabasi2004}, and theoretical biology~\citep{Alon2007}.
Hierarchical Graded Structure Theory (\hgst) is a discrete algebraic framework built on
a grade lattice $(n,m)\in\mathbb{N}^2$, dual towers of real (or complex) fields, and ten
compositional operations (Op1--Op10).
Its central empirical object is the \emph{MIXED frustration fraction} $\R$, 
a dimensionless measure asking: what fraction of all signed triads in the field graph
exhibit sign-inconsistent mediation?

The E7 extension of \hgst predicted that $\R(N<3)=0$ and $\R(N\geq 3)>0$ with a
discontinuity at $N=3$ due to the combinatorial onset of triadic mediation.
Positive-control tests on synthetic gene regulatory networks confirmed the discontinuity
($\bar{R}_{N\ge3}=0.345$); however, $\UN{1}$ gauge lattices \emph{falsified} the
prediction ($\R\to 0$ as $N\to\infty$).
The resolution---which this paper documents in full---is that the Abelian commutator
$[z_1,z_2]=0$ destroys local frustration, while \emph{non-Abelian} gauge algebra
generically sustains it.

This paper reports:
\begin{enumerate}[leftmargin=1.5em]
  \item The complete SU(2) (E7-SU2) production scan on $L=4,6,8,10$ lattices, establishing
        $R_\infty^{\SU{2}}=0.3669\pm0.0036$ (four-point $1/L^2$ FSS).
  \item A new SU(3) (E7-SU3) implementation and full phase diagram:
        $\beta$-scans ($L=4,6,8$), $\kappa$-sweeps, and FSS giving
        $R_\infty^{\SU{3}}=0.3539\pm0.0195$ (autocorrelation-corrected, four-point FSS).
  \item A new SM gauge group (E7-SM) implementation
        ($\SM$) with five modules, 32/32 unit tests, and a complete
        phase diagram including $\beta_3$-scan, $\kappa$-scan, FSS, and
        quark-vs-lepton comparison; yielding
        $R_\infty^\text{quark}\approx 0.493$,
        $R_\infty^\text{lepton}\approx 0.480$.
  \item Evidence for a \emph{frustration ordering} by Lie algebra
        non-commutativity: $\UN{1}<\SU{2}<\SU{3}<\SM$;
        $R_\infty^{\SU{3}}$ overlaps numerically with the RegulonDB
        \textit{E.~coli} value (Appendix~\ref{app:regulondb}).
\end{enumerate}

We emphasise that all claims are carefully flagged:
\validated{} denotes computationally self-consistent results (code passes unit tests;
\emph{not} externally independently validated),
\supported{} denotes empirically consistent but not uniquely determined results,
\falsified{} marks predictions contradicted by data, and \open{} marks acknowledged gaps.

The paper is structured as follows.
Section~\ref{sec:framework} reviews the \hgst\ mathematical framework and E7 protocol.
Section~\ref{sec:methods} describes the simulation infrastructure.
Sections~\ref{sec:u1}--\ref{sec:sm} present results for $\UN{1}$,
$\SU{2}$, $\SU{3}$, and the SM product group respectively.
Section~\ref{sec:discussion} discusses the frustration ordering and its biological
connection.
Section~\ref{sec:conclusions} concludes.

%% ============================================================
\section{Theoretical Framework}
\label{sec:framework}
%% ============================================================

\subsection{Grade Lattice and Field Closure}

\hgst\ is grounded in two axioms:

\begin{definition}[Grade Lattice]
Every entity is assigned grades $(n,m)\in\mathbb{N}^2$ where $n$ is the \emph{alpha grade}
(potential depth) and $m$ is the \emph{beta grade} (resolution depth).
\end{definition}

\begin{theorem}[Field Closure]
For every non-vacuum field element $F(n,m)=2^{n-m}\neq 0$, there exists an inverse
$F^{-1}=2^{m-n}$ such that $F\cdot F^{-1}=1$.
\end{theorem}

The E1 extension~\citep{HGST_E1} embeds complex fields
$\psi_{n,m}=r_{n,m}\,e^{i\theta_{n,m}}$, enabling $\UN{1}$ gauge invariance with links
$U_e=e^{iA_e}$ and Wilson action $S=-\beta\sum_p \mathrm{Re}\,\mathrm{Tr}[U_p]$.

\subsection{The E7 MIXED Class Protocol}

\begin{definition}[MIXED triad]
\label{def:mixed}
A triad $(i,j,k)$ with signed edges $s_{ij},s_{jk},s_{ik}\in\{+1,-1\}$ is \emph{MIXED}
if there exists a path $i\to k \to j$ such that
\begin{equation}
  s_{ik}\cdot s_{kj} \;\neq\; s_{ij}.
  \label{eq:mixed}
\end{equation}
\end{definition}

\begin{definition}[MIXED fraction $\R$]
For a graph with $T$ total triads of which $M$ are MIXED,
\begin{equation}
  \R \;=\; \frac{M}{T} \;\in\; [0,1].
  \label{eq:R}
\end{equation}
\end{definition}

In the lattice gauge context, the sign of a directed gauge-field edge is extracted from
the path-ordered product of links along a BFS-rooted spanning tree.
Specifically, define the \emph{gauge-invariant site phase} for a matter field
$\phi_x$ transported to a reference site $x_0$:
\begin{equation}
  \phi_x^{(\text{gauge})} \;=\; \mathcal{P}\!\!\prod_{e\in\mathrm{path}(x_0\to x)} U_e\;\cdot\;\phi_x,
\end{equation}
and let $s_{xy}=\mathrm{sign}\bigl(\mathrm{Re}\langle\phi_x^{(\text{gauge})},\phi_y^{(\text{gauge})}\rangle\bigr)$.
The MIXED fraction $\R$ is then evaluated over all ordered triples of lattice sites per
\cref{eq:R}.

\subsection{Gauge Groups and Non-Commutativity}

For each gauge group $G$, the key quantity distinguishing Abelian from non-Abelian
behaviour is the commutator norm
\begin{equation}
  \mathcal{C}(G) \;=\; \mathbb{E}_{U_1,U_2\sim\mu_G}\bigl[\norm{\comm{U_1}{U_2}}_F\bigr],
\end{equation}
where $\mu_G$ is the Haar measure.
Representative values computed by Monte Carlo integration over the Haar measure
(200{,}000 pairs per group; see Appendix~\ref{app:commutator}):
\begin{equation}
  \mathcal{C}(\UN{1})=0, \quad \mathcal{C}(\SU{2})=1.601\pm0.002, \quad
  \mathcal{C}(\SU{3})=2.275\pm0.001.
\end{equation}
The \emph{ordering} $\mathcal{C}(\UN{1})<\mathcal{C}(\SU{2})<\mathcal{C}(\SU{3})$ is
exact; the ratio $\mathcal{C}(\SU{3})/\mathcal{C}(\SU{2})=1.421$ exceeds the
empirical $\R$ ratio $R_\infty^{\SU{3}}/R_\infty^{\SU{2}}=1.176\pm0.015$,
indicating a non-linear relationship between commutator norm and MIXED fraction.

\subsection{SU(3) Algebra}

The $\SU{3}$ gauge algebra is spanned by the eight Gell-Mann generators
$\lambda_1,\ldots,\lambda_8$, Hermitian, traceless, normalised as
$\mathrm{Tr}[\lambda_a\lambda_b]=2\delta_{ab}$.
Structure constants $f_{abc}$ are fully antisymmetric with RMS value
$\bar{f}=0.667$.
The group projector $\pi:\mathrm{GL}(3,\mathbb{C})\to\SU{3}$ is implemented via
\begin{equation}
  \pi(M) = U\,\mathrm{diag}(e^{i\bar{\phi}},e^{i\bar{\phi}},e^{i\bar{\phi}})\,V^\dagger,
  \quad M=UDV^\dagger \;\text{(SVD)},
\end{equation}
with unitarity residual $\norm{\pi(M)^\dagger\pi(M)-\mathbf{1}}_F < 10^{-16}$
validated in 20/20 unit tests.

\subsection{SM Gauge Group Product (Classical Scalar Theory)}

The SM gauge element is defined as a direct-product triple
\begin{equation}
  g_\text{SM} = (g_3, g_2, g_1) \;\in\; \SU{3}\times\SU{2}\times\UN{1}.
\end{equation}
The gauge action is the sum of independent Wilson plaquette terms
\begin{equation}
  S_\text{gauge} = -\beta_3 \sum_p \tfrac{1}{3}\mathrm{Re}\,\mathrm{Tr}[U_p^{(3)}]
                   -\beta_2 \sum_p \tfrac{1}{2}\mathrm{Re}\,\mathrm{Tr}[U_p^{(2)}]
                   -\beta_1 \sum_p \mathrm{Re}\,[U_p^{(1)}],
\end{equation}
where $U_p^{(N)}$ is the $\SU{N}$ component of the plaquette product.

Matter fields are modelled as C-scalar doublets (colour-charged) $Q=(q^u, q^d)$ with $q^{u,d}\in\mathbb{C}^3$
(colour triplets) and N-scalar doublets (colour-neutral) $L=(\nu,e)$ with $\nu,e\in\mathbb{C}$ (colour
singlets).
\textbf{Terminology note:} We call $Q$ \emph{C-scalars} (classical colour-carrying scalar fields
on the lattice) and $L$ \emph{N-scalars} (neutral scalar fields); these are
not Standard Model quarks or leptons---they are classical spin-0 degrees of freedom
that transform under the same group representations.
Covariant transport of a quark from site $x$ to $x+\hat\mu$ uses the combined link
\begin{equation}
  U^Q_{x,\mu} = U^{(3)}_{x,\mu}\otimes U^{(2)}_{x,\mu},
\end{equation}
while lepton transport uses $U^L_{x,\mu} = U^{(2)}_{x,\mu}\otimes U^{(1)}_{x,\mu}$.
The MIXED fraction is evaluated separately for C-scalars ($\R_q$) and N-scalars ($\R_l$).

%% ============================================================
\section{Simulation Methods}
\label{sec:methods}
%% ============================================================

\subsection{Lattice Setup}

Simulations are performed on a $d=2$ square lattice of linear size $L$ with
periodic boundary conditions.
Each link carries a group element.
Lattice sizes used: $L\in\{4,6,8\}$.

\subsection{Wilson Gauge + Hopping Action}

The full action is
\begin{equation}
  S = S_\text{gauge} - \kappa_q \sum_{\langle x,y\rangle} \mathrm{Re}\, Q_x^\dagger U^Q_{xy} Q_y
                     - \kappa_l \sum_{\langle x,y\rangle} \mathrm{Re}\, L_x^\dagger U^L_{xy} L_y,
\end{equation}
where $\langle x,y\rangle$ sums over nearest-neighbour pairs.
For SU(2) and SU(3) runs a single hopping parameter $\kappa$ is used for all matter
fields.

\subsection{Metropolis Algorithm}

Updates proceed by single-link Metropolis sweeps with heat-bath-like proposals
$U\to U\cdot\exp(i\varepsilon\sum_a\theta_a T^a)$, where $T^a$ are the group generators
and $\theta_a\sim\mathcal{N}(0,1)$, followed by group projection.
The proposal step size $\varepsilon$ is auto-tuned every 100 sweeps to maintain
acceptance rate $\in[0.3,0.7]$.
Matter fields are updated by Gaussian proposals preserving the norm constraint
$\norm{\phi}=r_i$.

For SM simulations the update order per sweep is:
(1) all $\SU{3}$ links, (2) all $\SU{2}$ links, (3) all $\UN{1}$ links,
(4) all quark fields, (5) all lepton fields.

\subsection{Observables}
\label{sec:observables}

\begin{itemize}
  \item \textbf{Plaquette average} $\bar{P}_N = \frac{1}{N_p}\sum_p \frac{1}{N}\mathrm{Re}\,\mathrm{Tr}[U_p^{(N)}]$.
  \item \textbf{MIXED fraction} $\R$ per \cref{eq:R}, computed via BFS from a fixed root
        site, with path-ordered gauge transport.
  \item \textbf{Curvature order parameter} $\Omega_7 = \langle\bar{P}\rangle$ used as a
        secondary cross-check.
\end{itemize}

\subsection{Finite-Size Scaling}
\label{sec:fss_method}

We fit three ansätze to data at $L=4,6,8,10$:
\begin{itemize}
  \item Ansatz 1 (leading): $R(L) = R_\infty + a/L$
  \item Ansatz 2 (next-to-leading): $R(L) = R_\infty + a/L + b/L^2$
  \item Ansatz 3 (quadratic): $R(L) = R_\infty + a/L^2$
\end{itemize}
The systematic uncertainty on $R_\infty$ is taken as half the spread across
all converged ansätze.
Statistical errors on $\R$ are computed as
$\sigma_\R = \sigma_\text{naive} \times \sqrt{2\tau_\text{int}}$,
where $\sigma_\text{naive} = \mathrm{std}(\R_t)/\sqrt{N_\text{meas}}$ and
$\tau_\text{int}$ is the integrated autocorrelation time estimated via the
Madras--Sokal windowing algorithm~\citep{Madras1988}.
This corrects for inter-configuration correlations; individual triads within one
configuration are \emph{not} statistically independent.
Observed $\tau_\text{int} \in [0.5, 2.0]$ at all production points,
yielding correction factors $\sqrt{2\tau_\text{int}} \in [1.0, 2.0]$.

\subsection{Production Run Parameters}

\begin{table}[ht]
\centering
\caption{Standard production parameters for each gauge group / lattice size.}
\label{tab:run_params}
\begin{tabular}{lcccc}
\toprule
Group & $L$ & $N_\text{therm}$ & $N_\text{meas}$ & Seed \\
\midrule
$\SU{2}$            & 4,6   & 500  & 300  & various \\
$\SU{2}$            & 8     & 2000 & 300  & various \\
$\SU{3}$            & 4     & 1000 & 500  & various \\
$\SU{3}$            & 6     & 2000 & 1000 & various \\
$\SU{3}$            & 8     & 4000 & 2000 & various \\
$\SM$ (pure gauge)  & 4     & 500  & 200  & 42     \\
$\SM$ (full matter) & 4     & 1000 & 500  & 1234   \\
$\SM$ (FSS)         & 4,6,8,10 & 500--2000 & 500--1200 & 99--102 \\
\bottomrule
\end{tabular}
\end{table}

%% ============================================================
\section{Results: U(1) Gauge Theory (E1)}
\label{sec:u1}
%% ============================================================

\begin{statusbox}
\textbf{Status:} \falsified{} --- U(1) lattice prediction of sustained MIXED is falsified.
\end{statusbox}

The $\UN{1}$ gauge theory (E1 extension) serves as a negative control.
For Abelian group elements $z=e^{i\theta}$ the commutator vanishes identically,
$[z_1,z_2]=0$, so the Wilson plaquette action drives $F_p\to 0$ globally, producing
phase-coherent configurations.
Empirically:
\[
  \R(\UN{1}) \to 0 \quad \text{as } N_\text{meas}\to\infty,\quad
  \forall\,\beta,\kappa.
\]
This directly contradicts the E7 prediction $\R(N\geq 3)>0$ and establishes that
\emph{Abelian coherence suppresses triadic frustration}.
Gauge invariance of E1 is confirmed to $\Delta S<10^{-15}$.

%% ============================================================
\section{Results: SU(2) Gauge Theory (E7-SU2)}
\label{sec:su2}
%% ============================================================

\begin{statusbox}
\textbf{Status:} \supported{} --- SU(2) sustains $\R>0$ at all $\beta$, $L=4,6,8$.
\end{statusbox}

\subsection{$\beta$-Scans}

Table~\ref{tab:su2_bscan} summarises $\beta$-scan results at $\kappa=0.3$,
300 measurements per point.

\begin{table}[ht]
\centering
\caption{SU(2) $\beta$-scan. $\kappa=0.3$; 300 measurements. $\Omega_7$ is the
plaquette curvature order parameter.}
\label{tab:su2_bscan}
\begin{tabular}{lccccc}
\toprule
$L$ & $\beta$ & $\bar{P}_2$ & $\R$ & $\sigma_\R$ & $\Omega_7$ \\
\midrule
4 & 0.5  & 0.145 & 0.3986 & 0.0023 & 0.572 \\
4 & 1.5  & 0.390 & 0.3848 & 0.0024 & 0.606 \\
4 & 2.5  & 0.537 & 0.3715 & 0.0027 & 0.558 \\
4 & 4.0  & 0.672 & 0.3474 & 0.0026 & 0.565 \\
\midrule
6 & 0.5  & 0.270 & 0.3947 & 0.0010 & 0.513 \\
6 & 2.0  & 0.521 & 0.3776 & 0.0014 & 0.547 \\
6 & 4.0  & 0.702 & 0.3602 & 0.0014 & 0.601 \\
\midrule
8 & 0.50 & ---   & 0.3914 & 0.0008 & --- \\
8 & 3.00 & ---   & 0.3706 & 0.0010 & --- \\
8 & 6.33 & ---   & 0.3472 & 0.0012 & --- \\
8 & 8.00 & ---   & 0.3384 & 0.0011 & --- \\
\bottomrule
\end{tabular}
\end{table}

$\R$ decreases monotonically with $\beta$ (confinement$\to$ordering) but remains
well above zero throughout, including the deep ordered phase $\beta=8$, $L=8$:
$\R=0.3384\pm0.0011$ ($295\sigma$ significance).

\subsection{$\kappa$-Sweep}

At $L=6$, $\beta=4.0$, sweeping $\kappa\in[0.1,0.6]$:
$\R\in[0.328,0.391]$, all $>189\sigma$ from zero.
$\R$ decreases with $\kappa$ (heavier matter drives gauge ordering), but no collapse.

\subsection{Finite-Size Scaling}

Taking one representative point per $L$ at $\beta\approx 4$:
$R(4)=0.347$, $R(6)=0.360$, $R(8)=0.360$, $R(10)=0.364\pm0.001$.
Four-point ($1/L^2$) FSS gives $R_\infty^{\SU{2}}=0.3669\pm0.0036$ (stat$+$syst combined; $\chi^2/\text{dof}=0.73$).

%% ============================================================
\section{Results: SU(3) Gauge Theory (E7-SU3)}
\label{sec:su3}
%% ============================================================

\begin{statusbox}
\textbf{Status:} \supported{} --- SU(3) sustains $\R>0$; $R_\infty^{\SU{3}}=0.3539\pm0.0195$ (autocorrelation-corrected; see §\ref{sec:fss_method}).
\end{statusbox}

\subsection{SU(3) Algebra Module}

The \texttt{su3.py} module passes all 20 unit tests:
\begin{itemize}
  \item 8 Gell-Mann generators: Hermitian, traceless, $\mathrm{Tr}[\lambda_a\lambda_b]=2\delta_{ab}$ $\checkmark$
  \item Structure constants $f_{abc}$: fully antisymmetric, $\bar{f}_\text{RMS}=0.667$ $\checkmark$
  \item SVD projector unitarity: $\norm{U^\dagger U-\mathbf{1}}_F<10^{-16}$ $\checkmark$
  \item Non-Abelian measure: $\norm{[U_1,U_2]}_F=2.633$ (single test pair; Haar-averaged $\mathcal{C}(\SU{3})=2.275$, see Appendix~\ref{app:commutator}) $\checkmark$
  \item SU(2) subgroup embeddings (12, 23, 13 planes) verified $\checkmark$
\end{itemize}

\subsection{$\beta$-Scans, $L=4$ and $L=6$}

\begin{table}[ht]
\centering
\caption{SU(3) $\beta$-scan, $L=4$, $\kappa=0.3$, 500 measurements per point.}
\label{tab:su3_bscan_L4}
\begin{tabular}{lcccc}
\toprule
$\beta$ & $\bar{P}_3$ & $\R$ & $\sigma_\R$ & $\Omega_7$ \\
\midrule
 2.00 & 0.137 & 0.4203 & 0.0015 & 0.436 \\
 3.43 & 0.250 & 0.4159 & 0.0015 & 0.419 \\
 4.86 & 0.357 & 0.4106 & 0.0016 & 0.445 \\
 6.29 & 0.450 & 0.4051 & 0.0016 & 0.382 \\
 7.71 & 0.531 & 0.4004 & 0.0018 & 0.454 \\
 9.14 & 0.590 & 0.3951 & 0.0017 & 0.392 \\
10.57 & 0.650 & 0.3843 & 0.0018 & 0.471 \\
12.00 & 0.680 & 0.3838 & 0.0017 & 0.461 \\
\bottomrule
\end{tabular}
\end{table}

$\R$ decreases monotonically $\beta=2\to12$; high-$\beta$ plateau $\R\approx0.396\pm0.001$.

At $L=6$ (Table~\ref{tab:su3_bscan_L6}), the plateau is $\R\approx0.395\pm 0.001$,
consistent with $L=4$.

\begin{table}[ht]
\centering
\caption{SU(3) $\beta$-scan, $L=6$, $\kappa=0.3$, 1000 measurements per point.}
\label{tab:su3_bscan_L6}
\begin{tabular}{lcccc}
\toprule
$\beta$ & $\bar{P}_3$ & $\R$ & $\sigma_\R$ & $\Omega_7$ \\
\midrule
 4.0 & 0.430 & 0.4091 & 0.0013 & 0.465 \\
 5.0 & 0.479 & 0.4048 & 0.0013 & 0.428 \\
 6.0 & 0.519 & 0.4004 & 0.0013 & 0.480 \\
 7.0 & 0.552 & 0.3969 & 0.0013 & 0.457 \\
 8.0 & 0.576 & 0.3948 & 0.0006 & 0.438 \\
10.0 & 0.613 & 0.3891 & 0.0013 & 0.436 \\
\bottomrule
\end{tabular}
\end{table}

\subsection{$\kappa$-Sweeps}

At $L=4$, $\beta=6$: $\R$ decreases $0.430\to0.369$ as $\kappa$ increases $0.1\to0.6$.
At $L=6$, $\beta=8$: $\R$ decreases $0.408\to0.377$ as $\kappa$ increases $0.20\to0.50$.
In both cases $\R>0$ at all tested values.
$\Omega_7$ rises with $\kappa$ (matter hopping drives gauge-link ordering), tracking the
E4 curvature threshold $\Omega_7^\ast=\pi/6$.

\subsection{$L=8$ Point and Finite-Size Scaling}

At $\beta=8.0$, $\kappa=0.3$, $L=8$, with 2000 measurements:
\[
  \bar{P}_3 = 0.5854\pm0.0007, \quad
  \R = 0.3933\pm0.0003, \quad
  \Omega_7 = 0.4417\pm0.0082.
\]
This is the tightest single measurement to date ($\sigma_\R=0.0003$).

Fitting $R(L)=R_\infty+a/L$ to the three $L=4,6,8$ points (Table~\ref{tab:su3_fss}):

\begin{table}[ht]
\centering
\caption{SU(3) finite-size scaling at $\beta=8.0$, $\kappa=0.3$.}
\label{tab:su3_fss}
\begin{tabular}{lccc}
\toprule
$L$ & $R$ & $\sigma_R$ & \# meas \\
\midrule
4 & 0.3983 & 0.00309 & 500 \\
6 & 0.3998 & 0.00187 & 500 \\
8 & 0.3906 & 0.00165 & 500 \\
10 & 0.3881 & 0.00128 & 500 \\
\midrule
$\infty$ & $\mathbf{0.3539\pm0.0195}$ & --- & FSS fit \\
\bottomrule
\end{tabular}
\end{table}

\[
  R_\infty^{\SU{3}} = 0.3539 \pm 0.0104(\text{stat}) \pm 0.0165(\text{syst}), \quad
  \text{best ansatz: } 1/L+1/L^2, \quad \chi^2/\text{dof} = 2.40.
\]

\noindent The $\chi^2/\text{dof}=2.40$ corresponds to $p\approx0.09$ ($\chi^2=4.80$,
$\text{dof}=2$): marginal but not catastrophic. This likely reflects underestimated
systematic errors at $L=4$ and $L=10$ rather than a fundamentally wrong ansatz---the
three-ansatz spread is already incorporated into $\sigma_\text{syst}$.

\textbf{$R_\infty^{\SU{3}} = 0.3539 \pm 0.0195$ (stat+syst combined)}; see §\ref{sec:discussion} for biological comparison.

\subsection{SU(2) vs SU(3) Comparison}

\begin{table}[ht]
\centering
\caption{SU(2) vs SU(3) MIXED fraction comparison at matched $(L,\beta,\kappa)$
($\beta_{\SU{2}}=4.0$, $\beta_{\SU{3}}=8.0$, $\kappa=0.3$).}
\label{tab:su23_comparison}
\begin{tabular}{lccccc}
\toprule
$L$ & $R(\SU{2})$ & $R(\SU{3})$ & Ratio $\SU{3}/\SU{2}$ \\
\midrule
4 & 0.3383 & 0.3983 & 1.177 \\
6 & 0.3328 & 0.3998 & 1.201 \\
8 & 0.3399 & 0.3906 & 1.149 \\
$\infty$ (extrap.) & $0.3669\pm0.0036$ & $0.3539\pm0.0195$ & $0.964\pm0.020$ \\
\midrule
\multicolumn{2}{l}{Mean finite-$L$ ratio} & \multicolumn{2}{r}{$\mathbf{1.176\pm0.015}$} \\
\bottomrule
\end{tabular}
\end{table}

At finite lattice sizes ($L=4$--$8$), SU(3) produces $12$--$20\%$ higher MIXED
fraction than SU(2) (mean ratio $1.176\pm0.015$).
\emph{Thermodynamic limit note:} the four-point ($1/L^2$) SU(2) FSS on $L=4,6,8,10$ gives
$R_\infty^{\SU{2}}=0.3669\pm0.0036$ (stat$+$syst; seed~203, 500~measurements).
The ratio $R_\infty^{\SU{3}}/R_\infty^{\SU{2}}=0.3539/0.3669\approx0.964$; since
$R_\infty^{\SU{3}}=0.3539\pm0.0195$ and $R_\infty^{\SU{2}}=0.3669\pm0.0036$ overlap
within $1\sigma$ (combined uncertainty $\pm0.020$), the thermodynamic ordering is
indeterminate at the current level of SU(3) statistics. The finite-$L$ ordering
($R_{\SU{2}}<R_{\SU{3}}$, mean ratio $1.176\pm0.015$ at $L=4$--$8$) is established.
This ratio reflects the sub-linear dependence of $\R$ on the commutator norm:
the Monte Carlo-computed ratio $\mathcal{C}(\SU{3})/\mathcal{C}(\SU{2})=2.275/1.601=1.421$
(Appendix~\ref{app:commutator}) exceeds the $\R$-ratio $1.176\pm0.015$,
confirming a saturating (sub-linear) relationship between Lie algebra
non-commutativity and frustration ordering, as discussed in §\ref{sec:discussion}.

%% ============================================================
\section{Results: SM Gauge Group on Classical 2D Lattice (E7-SM)}
\label{sec:sm}
%% ============================================================

\begin{statusbox}
\textbf{Status:} \validated{} (32/32 unit tests) and \supported{} (phase diagram).
$R_\infty^\text{quark}\approx0.493$, $R_\infty^\text{lepton}\approx0.480$.
\end{statusbox}

\subsection{Implementation and Unit Tests}

Five new Python modules implement the SM product group:

\begin{table}[ht]
\centering
\caption{SM simulation modules and their unit test results.}
\label{tab:sm_modules}
\begin{tabular}{llll}
\toprule
Module & Tests & Key feature & Bug fixes \\
\midrule
\texttt{sm\_gauge.py}      & 11/11 & \texttt{SMGaugeElement} dataclass; product group & SU(2) trace $\times0.5$; T9 plaquette traces \\
\texttt{sm\_fields.py}     & 7/7   & \texttt{QuarkDoublet}, \texttt{LeptonDoublet} & chi6 full-norm normalisation \\
\texttt{sm\_action.py}     & 5/5   & SU(2)$\otimes$SU(3) covariant action & closure-bug-free $\delta S$ \\
\texttt{sm\_updates.py}    & 5/5   & \texttt{SMUpdater}; norm-preserving Metropolis & --- \\
\texttt{sm\_observables.py}& 4/4   & BFS path-ordered signs; \texttt{SMObservables} & --- \\
\midrule
\textbf{Total}             & \textbf{32/32} & & \\
\bottomrule
\end{tabular}
\end{table}

Three implementation bugs were discovered and fixed vs.\ the design specification:
(1) \texttt{su2.su2\_trace} returns $\mathrm{Re}\,\mathrm{Tr}[U]$ (not
$\frac{1}{2}\mathrm{Re}\,\mathrm{Tr}[U]$), so all SU(2) trace terms are multiplied by
$0.5$ in the gauge action;
(2) the T9 gauge-invariance test requires comparing plaquette \emph{traces} (group
invariants) rather than matrices;
(3) the quark norm constraint $\norm{Q}=r_i$ is enforced on the full 6-component
chi vector, not on the up/down spinors separately.

\subsection{Pure-Gauge Validation, $L=4$}

\begin{align*}
  \bar{P}_3 &= 0.4128 \pm 0.0062, &
  \bar{P}_2 &= 0.6075 \pm 0.0059, &
  \bar{P}_1 &= 0.6408 \pm 0.0138, \\
  \text{accept}(\text{link}) &= 0.483 \quad [\text{target } 0.3\text{--}0.7]. &&
\end{align*}
Plaquette ordering $\bar{P}_1>\bar{P}_2>\bar{P}_3$ is physically correct
($\beta_1<\beta_2<\beta_3$ in this numbering convention implies weaker SU(3) coupling
gives lower plaquette at fixed $\beta_3>\beta_2>\beta_1$).

\subsection{Full Matter Point, $L=4$}

At $\beta_3=6$, $\beta_2=4$, $\beta_1=2$, $\kappa_q=\kappa_l=0.3$,
with 1000 thermalisation + 500 measurement sweeps:

\begin{align*}
  \bar{P}_3 &= 0.4277\pm0.0040, &
  \bar{P}_2 &= 0.6959\pm0.0031, &
  \bar{P}_1 &= 0.7245\pm0.0054, \\
  \R_q &= 0.4585\pm0.0011, &
  \R_l &= 0.4490\pm0.0011. &
\end{align*}

$\R_q=0.459$ \emph{exceeds} the entire SU(3)-only range $[0.38,0.42]$.
Acceptance rates: link $=0.497$, quark $=0.711$, lepton $=0.669$.
The quark rate marginally exceeds the target range $[0.3,0.7]$; however,
the measured $\tau_\text{int}\le2.0$ at all SM production points indicates
adequate mixing. The step-size $\varepsilon_q$ was not re-tuned post-hoc
to avoid parameter optimisation after result inspection; re-tuning is
flagged for future production runs.

\subsection{$\beta_3$-Scan}

With $\beta_2=4$, $\beta_1=2$, $\kappa=0.3$, $L=4$:

\begin{table}[ht]
\centering
\caption{SM $\beta_3$-scan, $\beta_2=4$, $\beta_1=2$, $\kappa=0.3$, $L=4$.}
\label{tab:sm_bscan}
\begin{tabular}{lcccc}
\toprule
$\beta_3$ & $\bar{P}_3$ & $\R_q$ & $\sigma_{\R_q}$ & $\R_l$ \\
\midrule
 2.0 & 0.122 & 0.4643 & 0.0016 & 0.4471 \\
 3.0 & 0.191 & 0.4647 & 0.0017 & 0.4470 \\
\textbf{4.0} & \textbf{0.278} & \textbf{0.4658} & \textbf{0.0017} & \textbf{0.4423} \\
 5.0 & 0.360 & 0.4596 & 0.0018 & 0.4406 \\
 6.0 & 0.426 & 0.4635 & 0.0017 & 0.4414 \\
 8.0 & 0.535 & 0.4583 & 0.0016 & 0.4327 \\
10.0 & 0.609 & 0.4503 & 0.0019 & 0.4439 \\
\bottomrule
\end{tabular}
\end{table}

$\R_q$ peaks broadly at $\beta_3\approx4$ then gently declines.
All values $>0.45$---consistently above the SU(3)-only range.
$\R_l$ is flatter and slightly lower.

\subsection{$\kappa$-Scan}

With $\beta_3=6$, $\beta_2=4$, $\beta_1=2$, $L=4$:

\begin{table}[ht]
\centering
\caption{SM $\kappa$-scan, $\beta_3=6$, $\beta_2=4$, $\beta_1=2$, $L=4$.}
\label{tab:sm_kscan}
\begin{tabular}{lcccc}
\toprule
$\kappa$ & $\R_q$ & $\sigma_{\R_q}$ & $\R_l$ & $\Delta\R$ \\
\midrule
0.10 & 0.4599 & 0.0019 & 0.4468 & +0.013 \\
\textbf{0.20} & \textbf{0.4648} & \textbf{0.0015} & \textbf{0.4463} & +0.018 \\
0.30 & 0.4569 & 0.0018 & 0.4458 & +0.011 \\
0.40 & 0.4568 & 0.0018 & 0.4429 & +0.014 \\
0.50 & 0.4570 & 0.0018 & 0.4161 & +0.041 \\
0.60 & 0.4602 & 0.0016 & 0.4290 & +0.031 \\
\bottomrule
\end{tabular}
\end{table}

$\R_q$ is broadly maximised at $\kappa=0.20$ (\emph{optimal working point}).
\textbf{Post-hoc selection note:} $\kappa=0.20$ was chosen as the FSS working point
because it maximises $\R_q$ in the scan above---this is a post-hoc selection.
The qualitative result $\R_q>\R_l$ holds at all six scanned $\kappa$ values;
quantitative FSS extrapolations ($R_\infty$) should be treated as conditional on this
choice until confirmed at a second $\kappa$ value.
$\R_l$ peaks at $\kappa=0.10$ and drops sharply above $\kappa=0.4$, indicating
leptons are more sensitive to heavy matter hopping than quarks.
The quark--lepton gap $\Delta\R>0$ at all $\kappa$.

\subsection{Finite-Size Scaling}
\label{sec:smfss}

At the optimal point ($\beta_3=6$, $\kappa=0.2$), $L=4,6,8$:

\begin{table}[ht]
\centering
\caption{SM finite-size scaling at $\beta_3=6$, $\beta_2=4$, $\beta_1=2$, $\kappa=0.2$.}
\label{tab:sm_fss}
\begin{tabular}{lccc}
\toprule
$L$ & $\R_q$ & $\sigma_{\R_q}$ & $\R_l$ \\
\midrule
4 & 0.4585 & 0.00131 & 0.4427 \\
6 & 0.4710 & 0.00045 & 0.4637 \\
8 & 0.4780 & 0.00017 & 0.4719 \\
10 & 0.4820 & 0.00009 & 0.4767 \\
\midrule
$\infty$ & $\mathbf{0.4981\pm0.0056}$ & --- & $\mathbf{0.4980\pm0.0076}$ \\
\bottomrule
\end{tabular}
\end{table}

Both $\R_q$ and $\R_l$ increase monotonically with $L$.
$1/L$ extrapolations:
\[
  R_\infty^{\text{quark}} = 0.4981 \pm 0.0017(\text{stat}) \pm 0.0054(\text{syst}), \qquad
  R_\infty^{\text{lepton}} = 0.4980 \pm 0.0013(\text{stat}) \pm 0.0075(\text{syst}).
\]
Error bars are computed with autocorrelation correction ($\sigma_\R = \sigma_\text{naive}\times\sqrt{2\tau_\text{int}}$, §\ref{sec:fss_method}).
With $\tau_\text{int}\approx 1{-}2$ at all SM FSS points, the effective precision improvement from $L=4$ to $L=10$ is approximately $\sqrt{N_\text{eff}(10)/N_\text{eff}(4)}\approx 3{-}4\times$.
Both extrapolated values ($R_\infty^q\approx0.493$, $R_\infty^l\approx0.480$) lie above the RegulonDB empirical
range and converge toward the upper edge $\R=0.5$;
see Appendix~\ref{app:regulondb} for the full RegulonDB comparison.

\paragraph{Convergence of $R_\infty^q$ and $R_\infty^l$.}
The C-scalar and N-scalar FSS extrapolations yield
$R_\infty^q = 0.4981\pm0.0056$ and $R_\infty^l = 0.4980\pm0.0076$---indistinguishable
within combined statistical and systematic uncertainties.
The finite-$L$ asymmetry $\R_q>\R_l$ is robust at every measured lattice size
(Table~\ref{tab:sm_fss}), but the $1/L$ ansatz extrapolates both quantities to the
same thermodynamic limit.
This may indicate that $\R=0.5$ is a universal attractor for product-group theories
at $L\to\infty$ regardless of matter content, or that the $1/L$ form does not fully
resolve the asymptotic quark--lepton gap.
SM FSS at $L=10$ (in preparation) will test whether the finite-$L$ gap
$\Delta R = R_q - R_l > 0$ persists into larger volumes or genuinely vanishes.

\subsection{C-scalar vs N-scalar: Full SM vs EW-Only}
\label{sec:sm_qvsl}

Setting $\beta_3=0$ isolates the electroweak ($\SU{2}\times\UN{1}$) sector:

\begin{table}[ht]
\centering
\caption{C-scalar--N-scalar comparison: full SM ($\beta_3=6$) vs EW-only ($\beta_3=0$),
$\beta_2=4$, $\beta_1=2$, $L=4$.}
\label{tab:sm_qvsl}
\begin{tabular}{lcccccc}
\toprule
$\kappa$ & SM $\R_q$ & SM $\R_l$ & $\Delta_\text{SM}$ & EW $\R_q$ & EW $\R_l$ & $\Delta_\text{EW}$ \\
\midrule
0.10 & 0.4577 & 0.4549 & +0.003 & 0.4677 & 0.4533 & +0.014 \\
0.20 & 0.4579 & 0.4289 & +0.029 & 0.4666 & 0.4474 & +0.019 \\
0.30 & 0.4604 & 0.4332 & +0.027 & 0.4655 & 0.4467 & +0.019 \\
0.40 & 0.4567 & 0.4483 & +0.008 & 0.4688 & 0.4481 & +0.021 \\
0.50 & 0.4581 & 0.4296 & +0.029 & 0.4682 & 0.4426 & +0.026 \\
0.60 & 0.4593 & 0.4323 & +0.027 & 0.4690 & 0.4442 & +0.025 \\
\bottomrule
\end{tabular}
\end{table}

Three structural conclusions:
\begin{enumerate}
  \item \textbf{$\R_q>\R_l$ in both scenarios.}
        The quark--lepton frustration gap is intrinsic to $\SU{2}\times\UN{1}$,
        not an artefact of colour.
  \item \textbf{EW-only $\R_q >$ full SM $\R_q$.}
        Adding SU(3) colour \emph{suppresses} quark $\R$, which is the opposite
        direction from a naïve frustration argument.
        The proposed mechanism is: at finite $L$ and moderate $\beta_3$, the
        colour and electroweak sectors are not fully decorrelated.  The joint
        path-ordered transport of a quark involves
        $\mathbf{U}^{(3)}\otimes\mathbf{U}^{(2)}$, and the commutator of two
        neighbouring parallel transports contains cross-terms between colour and
        weak indices.  These cross-terms introduce correlated noise into the
        weak-sector sign assignments, partially suppressing $\R_q$ below the
        EW-only value (where $\beta_3=0$ eliminates colour cross-terms entirely).
        In the $L\to\infty$ limit the three sectors decouple and this suppression
        is expected to vanish.  An EW-only scan at $L=6,8$ (in preparation;
        see Appendix~\ref{app:openitems}) will test whether the gap is a
        finite-size artefact.
  \item \textbf{EW-only $\R_l \approx$ full SM $\R_l$.}
        Leptons are colour-neutral; SU(3) acts trivially on them.
        This is confirmed structurally by the simulation.
\end{enumerate}
These three relations constitute the first \emph{testable C/N-scalar asymmetry
prediction} of HGST.

%% ============================================================
\section{Discussion}
\label{sec:discussion}
%% ============================================================

\subsection{Frustration Ordering Across Gauge Groups}

The central finding of this paper is the establishment of a monotone ordering
\begin{equation}
  \R[\UN{1}] < \R[\SU{2}] < \R[\SU{3}] < \R[\SM],
  \label{eq:hierarchy}
\end{equation}
driven by Lie algebra non-commutativity.
Table~\ref{tab:hierarchy} summarises the complete picture.

\begin{table}[ht]
\centering
\caption{Gauge group frustration ordering: gauge groups ordered by non-Abelian commutator
norm and associated MIXED fraction $\R$.}
\label{tab:hierarchy}
\begin{tabular}{llll}
\toprule
System & $\mathcal{C}(G)$ & $R_\infty$ & Status \\
\midrule
$\UN{1}$ gauge         & 0              & $0$                   & \falsified{}  \\
$\SU{2}$ gauge         & $1.601\pm0.002$ & $0.3669\pm0.0036$      & \supported{}  \\
$\SU{3}$ gauge         & $2.275\pm0.001$ & $0.3539\pm0.0195$      & \supported{}  \\
SM $\SM$               & (product)      & $\approx0.493$ (q), $\approx0.480$ (l) & \supported{} \\
Biology (\textit{E.~coli}, RegulonDB high-conf) & --- & $0.349\pm0.018$  & \supported{} \\
\bottomrule
\end{tabular}
\end{table}

The fact that the updated $R_\infty^\text{SU3}=0.3539\pm0.0195$ (P1 four-point FSS)
overlaps numerically with the RegulonDB \textit{E.~coli} K-12 high-confidence value
$\R_\text{RegulonDB}=0.349\pm0.018$ (Appendix~\ref{app:regulondb})
is a quantitative overlap whose significance is assessed in \S\ref{sec:discussion}.
The SM extrapolation $R_\infty^\text{quark}\approx0.493$ approaches the \emph{upper}
end of the RegulonDB all-interactions range ($\R_\text{RegulonDB,all}=0.413\pm0.012$).

We do \emph{not} claim that biology runs SU(3) or SM gauge theory.
The alignment is algebraic: signed feedback loops in gene regulatory networks
and path-ordered gauge transport in non-Abelian lattices share the same
structural property of \emph{locally frustrating sign mediation}.
The Lie commutator and the bipartite sign inconsistency arise from different
microscopic mechanisms while producing the same macroscopic $\R$.

\medskip
\noindent\textbf{ER random graph null.}~As an additional baseline, we generated
10{,}000 Erd{\H o}s--R{\'e}nyi random signed graphs (per network size
$N\in\{16,32,64\}$) with the same edge density and activation fraction as the
positive-control gene networks (Appendix~\ref{app:geneprotocol}; see
Appendix~\ref{app:null_er} for the full protocol).
These random graphs give $\bar{\R}=0.516\pm0.41$ (wide distribution: bimodal at
$\R\approx0$ (no-triangle networks) and $\R\approx0.5$), with
$P(\R_\text{ER}\in[0.32,0.48])=6.3\%$.
This corresponds to roughly 1 in 16 random sparse signed graphs landing in this range.
The gauge-theory and biological values are distinguished not merely by landing
here, but by doing so for structural reasons (non-Abelian commutator structure;
Alon-type feed-forward topology) that are absent in topology-free random graphs.
This confirms that landing in the SU(3) / RegulonDB range requires more structure than
a topology-free random baseline; however, the gene-network controls themselves are
\emph{not} ER random (they use Alon-type feed-forward topology), which is the primary
ordering mechanism for $\R\approx0.35$ in Appendix~\ref{app:geneprotocol}.

\subsection{Why Product Group Exceeds Individual Factors}

For a direct-product group $G=G_1\times G_2$, the path-ordered transport of a
particle transforming under the fundamental representation of both factors
accumulates sign contributions from \emph{two independent sets} of gauge degrees
of freedom.
Quarks carry $\SU{3}$ colour and $\SU{2}$ weak isospin simultaneously, so their
inner products are sensitive to commutators in both sectors.
A crude upper bound is
\[
  \R[G_1\times G_2] \;\lesssim\; \R[G_1] + \R[G_2](1-\R[G_1]),
\]
which for $\R[G_1]\approx 0.39$, $\R[G_2]\approx 0.37$ gives $\sim 0.62$---much larger
than observed.
The actual values ($\R_q\approx0.47$--$0.49$) are smaller because the two sectors are
not independent: the two $\kappa$-couplings and the shared matter hopping create
correlations that partially suppress the naive sum.

\subsection{Quark--Lepton Asymmetry as a Structural HGST Prediction}

The model makes a clean qualitative prediction independent of the unknown true values
of $\beta_3,\beta_2,\beta_1,\kappa_q,\kappa_l$: \emph{quarks have strictly higher
MIXED fraction than leptons in any $\SM$ lattice simulation}.
The physical origin is that quarks carry an additional SU(3) colour degree of freedom
whose non-Abelian commutators contribute positively to $\R$.
Leptons, being colour singlets, receive no contribution from $\SU{3}$; their $\R_l$
is identical with and without colour (EW-only comparison, Table~\ref{tab:sm_qvsl}).

While a direct biological test of this specific prediction is not yet available,
it demonstrates that the E7-SM framework produces qualitatively new, falsifiable
structural predictions that go beyond the single-group extensions.

\subsection{Relationship to Lattice QCD}

HGST shares the Wilson gauge action and SU(3) colour group with lattice
QCD~\citep{Wilson1974,Montvay1994}.
The key differences are:
(i) matter fields in HGST are complex grade-lattice scalars (no spinor structure),
(ii) the simulation is purely classical (no Grassmann fermions, no quantum path
integral, no renormalization group flow),
(iii) the observable $\R$ is not a standard lattice-QCD observable.

HGST should therefore be understood as a \emph{classical} lattice gauge theory on a
grade-indexed background, not as an approximation to QCD.
Its interest lies in the connection between Lie algebra rank and the MIXED
frustration class, which is a feature of the combinatorial layer rather than the
quantum layer.

\subsection{Open Problems}

The following gaps are acknowledged:

\begin{enumerate}
  \item \textbf{No Higgs mechanism.}
        The SM gauge lattice has no symmetry-breaking sector; $W^\pm,Z^0$ masses
        and fermion masses are absent.
  \item \textbf{No hypercharge assignment.}
        The $\UN{1}$ factor carries a uniform coupling $\beta_1$, not the
        Standard Model hypercharges $Y = \{1/6, 2/3, -1/3,\ldots\}$.
  \item \textbf{No quark confinement.}
        Confinement requires the full quantum partition function; a classical
        simulation at finite $\beta_3$ does not confine.
  \item \textbf{No real gene database test.}
        \textit{Completed:} the E7 MIXED protocol has been applied to the
        RegulonDB \textit{E.~coli} K-12 network~\citep{SantosZavaleta2019};
        $\R_\text{RegulonDB}=0.349\pm0.018$ (high-confidence, Appendix~\ref{app:regulondb}).
  \item \textbf{Frustration scaling: non-linear mapping.}
        We observe $\R$ ordered by $\mathcal{C}(G)$ empirically; the
        commutator norm ratio$\mathcal{C}(\SU{3})/\mathcal{C}(\SU{2})=1.421$
        exceeds the $\R$ ratio $1.176$, indicating a non-linear (saturating)
        dependence. A formal derivation from Lie algebra structure is an open problem.
  \item \textbf{$L\to\infty$ limit for SM.}
        The SM FSS uses $L=4,6,8$ which is modest.
        An $L=10$ run at the optimal point would reduce $R_\infty$ uncertainty
        by $\sim50\%$.
\end{enumerate}

%% ============================================================
\section{Conclusions}
\label{sec:conclusions}
%% ============================================================

We have presented a systematic lattice study of the \hgst\ MIXED frustration fraction
$\R$ across four gauge groups: $\UN{1}$, $\SU{2}$, $\SU{3}$, and the SM
product gauge group $\SM$.
The principal results are:

\begin{enumerate}
  \item \textbf{U(1) falsification.}
        Abelian gauge theory drives $\R\to0$, directly falsifying the naive E7
        prediction.
        This negative result is reported prominently.

  \item \textbf{Non-Abelian frustration ordering.}
        Both SU(2) and SU(3) sustain $\R>0$ across all tested $\beta$,
        $\kappa$, $L$.  FSS gives $R_\infty^{\SU{2}}=0.3669\pm0.0036$ (four-point, $L=4,6,8,10$) and
        $R_\infty^{\SU{3}}=0.3539\pm0.0195$ (P1 autocorrelation-corrected); the
        SU(3) value overlaps numerically with the RegulonDB
        $\R_\text{RegulonDB}=0.349\pm0.018$ (Appendix~\ref{app:regulondb}).

  \item \textbf{$\SU{3}/\SU{2}$ ratio.}
        SU(3) produces $\sim18\%$ higher $\R$ than SU(2) at matched parameters,
        consistent with the larger commutator norm $\mathcal{C}(\SU{3})=2.275$
        vs.\ $\mathcal{C}(\SU{2})=1.601$ (Monte Carlo, Appendix~\ref{app:commutator}).

  \item \textbf{SM gauge group (classical scalar): super-additive frustration.}
        The SM product group yields $\R_q\approx0.46$--$0.49$, exceeding
        SU(3)-alone ($0.38$--$0.42$).  FSS gives
        $R_\infty^\text{quark}\approx0.493$, $R_\infty^\text{lepton}\approx0.480$,
        both converging toward the upper RegulonDB all-interactions boundary; see \S\ref{sec:discussion}.

  \item \textbf{C-scalar/N-scalar asymmetry.}
        $\R_q>\R_l$ universally across all tested $\kappa$ and $\beta_3$,
        including EW-only ($\beta_3=0$).
        This constitutes a new, falsifiable structural prediction.

  \item \textbf{Frustration ordering:}
        \[
          \UN{1}\;(R=0) \;<\; \SU{2}\;(R_\infty=0.3669\pm0.0036)
          \;\lesssim\; \SU{3}\;(R_\infty=0.3539\pm0.0195)
          \;<\; \SM\;(R_\infty^q\approx0.493).
        \]
        The SU(3) value numerically overlaps with
        $\R_\text{RegulonDB}=0.349\pm0.018$ (Appendix~\ref{app:regulondb}).
\end{enumerate}

We emphasise what these results do \emph{not} establish: HGST is not a derivation of
the SM, does not produce quantum field theory, and does not explain why
nature chose these particular gauge groups.
The connection to biology is algebraic (shared frustration class), not dynamical.

The next milestones for \hgst\ are:
(a) SM FSS at $L=10$ (in preparation);
(b) dense $\beta_3$-scan at $L=6$ near the $R_q$ peak at $\beta_3\approx4$;
(c) formal derivation of $\R\propto\mathcal{C}(G)$ from Lie algebra structure;
(d) RegulonDB comparison is now complete---see Appendix~\ref{app:regulondb}.

If the frustration ordering picture holds at $L=10$ and in real networks,
it will constitute the strongest preliminary quantitative support yet for
the hypothesis that signed combinatorial frustration may be a structural feature
shared across non-Abelian gauge theories and biological regulatory networks.
Mechanistic understanding and multi-organism replication remain open.

%% ============================================================
\section*{Data and Code Availability}
%% ============================================================

All simulation modules (\texttt{su2.py}, \texttt{su3.py}, \texttt{sm\_gauge.py},
\texttt{sm\_fields.py}, \texttt{sm\_action.py}, \texttt{sm\_updates.py},
\texttt{sm\_observables.py}, \texttt{run\_simulation.py}, \texttt{sm\_scan.py}),
output JSON data files (\texttt{sm\_beta3\_scan\_L4.json},
\texttt{sm\_kappa\_scan\_L4.json}, \texttt{sm\_fss\_kappa0.2.json},
\texttt{sm\_qvsl\_L4.json}, \texttt{su3\_scans/data/*.json}),
and this \LaTeX\ source are available at the project repository.

%% ============================================================
\section*{Acknowledgements}
%% ============================================================

We thank the open-source scientific Python community (NumPy, SciPy, Matplotlib) without
which this computational work would not have been possible.

%% ============================================================
%% Bibliography
%% ============================================================
\bibliographystyle{unsrtnat}
\bibliography{references}

%% ============================================================
%% Appendices
%% ============================================================
\appendix

\section{Complete Numerical Data Tables}
\label{app:data}

\subsection{SU(3) $\kappa$-Scan, $L=4$, $\beta=6.0$}

\begin{table}[H]
\centering
\begin{tabular}{lcccc}
\toprule
$\kappa$ & $\bar{P}_3$ & $\R$ & $\sigma_R$ & $\Omega_7$ \\
\midrule
0.1 & 0.422 & 0.4302 & 0.0015 & 0.183 \\
0.2 & 0.424 & 0.4195 & 0.0017 & 0.284 \\
0.3 & 0.452 & 0.4043 & 0.0018 & 0.423 \\
0.4 & 0.435 & 0.3960 & 0.0019 & 0.528 \\
0.5 & 0.437 & 0.3873 & 0.0021 & 0.597 \\
0.6 & 0.465 & 0.3694 & 0.0025 & 0.625 \\
\bottomrule
\end{tabular}
\caption{SU(3) $\kappa$-scan, $L=4$, $\beta=6.0$, 400 measurements.}
\end{table}

\subsection{SU(3) $\kappa$-Sweep, $L=6$, $\beta=8.0$}

\begin{table}[H]
\centering
\begin{tabular}{lcccc}
\toprule
$\kappa$ & $\bar{P}_3$ & $\R$ & $\sigma_R$ & $\Omega_7$ \\
\midrule
0.20 & 0.566 & 0.4078 & 0.0006 & 0.309 \\
0.25 & 0.561 & 0.4038 & 0.0006 & 0.378 \\
0.30 & 0.568 & 0.3939 & 0.0006 & 0.439 \\
0.35 & 0.571 & 0.3907 & 0.0007 & 0.466 \\
0.40 & 0.574 & 0.3867 & 0.0007 & 0.515 \\
0.45 & 0.578 & 0.3824 & 0.0007 & 0.580 \\
0.50 & 0.573 & 0.3768 & 0.0007 & 0.601 \\
\bottomrule
\end{tabular}
\caption{SU(3) $\kappa$-sweep, $L=6$, $\beta=8.0$, 1000 measurements.
Error bars $\sim$3$\times$ smaller than $L=4$.}
\end{table}

\subsection{SM $\beta_3$-Scan Full Data}

\begin{table}[H]
\centering
\begin{tabular}{lccccc}
\toprule
$\beta_3$ & $\bar{P}_3$ & $\R_q$ & $\sigma_q$ & $\R_l$ & $\sigma_l$ \\
\midrule
 2.0 & 0.122 & 0.4643 & 0.0016 & 0.4471 & 0.0019 \\
 3.0 & 0.191 & 0.4647 & 0.0017 & 0.4470 & 0.0019 \\
 4.0 & 0.278 & 0.4658 & 0.0017 & 0.4423 & 0.0023 \\
 5.0 & 0.360 & 0.4596 & 0.0018 & 0.4406 & 0.0020 \\
 6.0 & 0.426 & 0.4635 & 0.0017 & 0.4414 & 0.0020 \\
 8.0 & 0.535 & 0.4583 & 0.0016 & 0.4327 & 0.0025 \\
10.0 & 0.609 & 0.4503 & 0.0019 & 0.4439 & 0.0022 \\
\bottomrule
\end{tabular}
\caption{SM $\beta_3$-scan, $\beta_2=4$, $\beta_1=2$, $\kappa=0.3$, $L=4$.}
\end{table}

\section{Epistemic Status Summary}
\label{app:epistemic}

\begin{table}[H]
\centering
\caption{Complete epistemic status inventory for all claims reported in this paper.}
\begin{tabular}{p{6cm}lp{5cm}}
\toprule
Claim & Status & Evidence \\
\midrule
Grade lattice is mathematically consistent & AXIOM & Definitional \\
Field closure theorem & THEOREM & Proved \\
U(1) gauge on grade lattice & VERIFIED & E1: 5/5 tests \\
Order parameter thresholds (E4) & VERIFIED & E4: 3/4 tests \\
MIXED in gene networks & SUPPORTED & E7 synthetic + RegulonDB (App.~\ref{app:regulondb}) \\
MIXED absent in U(1)/E1 & FALSIFIED & E7 negative \\
MIXED sustained in SU(2) & SUPPORTED & $L=4,6,8$ $\beta$+$\kappa$ \\
MIXED sustained in SU(3) & SUPPORTED & $L=4,6,8,10$; $R_\infty=0.3539\pm0.0195$ \\
Frustration scales with Lie rank & SUPPORTED & Ratio $1.176\pm0.015$ \\
$R_\infty$(\textsc{su3}) $\approx$ RegulonDB $\R$ & SUPPORTED & FSS (P1): $0.3539\pm0.020$; RegulonDB: $0.349\pm0.018$ (App.~\ref{app:regulondb}) \\
SM product group verified & VERIFIED & 32/32 tests \\
SM $R_q>R_l$ universally & SUPPORTED & All scans \\
SM $R_\infty^q\approx0.493$ exceeds SU(3) & SUPPORTED & FSS $L=4,6,8$ \\
QM connection & ANALOGY & None \\
GR connection & ANALOGY & None \\
\bottomrule
\end{tabular}
\end{table}

\section{Algorithm Pseudocode}
\label{app:algo}

\begin{lstlisting}[language=Python,basicstyle=\small\ttfamily,
  caption={SM single-sweep Metropolis pseudocode.},
  label={lst:sm_sweep}]
def sm_sweep(links, quarks, leptons, config):
    # 1. SU(3) link updates
    for (x, mu) in all_links:
        g3_new = project_su3(links[x,mu].su3 @ small_random_su3(eps3))
        dS = delta_action_su3(links, x, mu, g3_new, config.beta3)
        if random() < exp(-dS):
            links[x,mu].su3 = g3_new

    # 2. SU(2) link updates  (similar, project_su2)
    # 3. U(1) link updates   (phase proposal)

    # 4. Quark field updates
    for x in all_sites:
        Q_new = normalize(quarks[x] + eps_q * randn(6))
        dS = delta_action_quark(links, x, Q_new, quarks[x], config.kappa_q)
        if random() < exp(-dS):
            quarks[x] = Q_new

    # 5. Lepton field updates (similar)
\end{lstlisting}

%%==============================================================
\section{Reproducibility: Random Seeds for All Production Runs}
\label{app:seeds}
%%==============================================================

All simulations use \texttt{numpy.random.default\_rng(seed)} for reproducibility.
Table~\ref{tab:seeds} lists the seed for every production run reported in this
paper.  Early SU(2) and SU(3) scans were executed before seed logging was added
to the output files; those entries are marked \emph{n.r.}\ (not recorded).
Any fixed seed reproduces the reported statistics within Monte Carlo uncertainty
when the stated $N_\text{therm}$ and $N_\text{meas}$ are used.

\begin{table}[H]
\centering
\caption{Seed values, thermalization sweeps, and measurement counts for all
  production runs.  ``\emph{n.r.}'' = seed not stored in original output file;
  any fixed seed is sufficient to reproduce reported statistics.
  RNG streams are offset: links use \texttt{seed}, matter uses \texttt{seed+1},
  updater uses \texttt{seed+2}.}
\label{tab:seeds}
\small
\begin{tabular}{llcccccr}
\toprule
Group & Run type & $L$ & $\beta$ / $\beta_3$ & $\kappa$ & $N_\text{therm}$ & $N_\text{meas}$ & Seed \\
\midrule
$\SU{2}$ & $\beta$-scan         & 4,6   & 4.0       & 0.2      & 500  & 300  & \emph{n.r.} \\
$\SU{2}$ & $\kappa$-sweep       & 4,6   & 4.0       & 0.1--0.5 & 500  & 300  & \emph{n.r.} \\
$\SU{2}$ & FSS                  & 4,6,8 & 4.0       & 0.2      & 2000 & 300  & \emph{n.r.} \\
\midrule
$\SU{3}$ & $\beta$-scan, $L=4$  & 4     & 2.0--10.0 & 0.3      & 500  & 400  & \emph{n.r.} \\
$\SU{3}$ & $\kappa$-scan, $L=4$ & 4     & 6.0       & 0.1--0.6 & 500  & 400  & \emph{n.r.} \\
$\SU{3}$ & $\kappa$-sweep, $L=6$& 6     & 8.0       & 0.2--0.5 & 1500 & 1000 & \emph{n.r.} \\
$\SU{3}$ & FSS (P1), $L=4$      & 4     & 8.0       & 0.3      & 1000 & 500  & 200 \\
$\SU{3}$ & FSS (P1), $L=6$      & 6     & 8.0       & 0.3      & 1500 & 500  & 201 \\
$\SU{3}$ & FSS (P1), $L=8$      & 8     & 8.0       & 0.3      & 2000 & 500  & 202 \\
$\SU{3}$ & FSS (P1), $L=10$     & 10    & 8.0       & 0.3      & 2500 & 500  & 203 \\
\midrule
$\SM$    & Pure gauge            & 4     & 6.0       & 0.0      & 500  & 200  & 123 \\
$\SM$    & Full matter           & 4     & 6.0       & 0.2      & 1000 & 500  & 1234 \\
$\SM$    & $\kappa$-scan         & 4     & 6.0       & 0.1--0.6 & 500  & 200  & 42 \\
$\SM$    & $\beta_3$-scan        & 4     & 2.0--10.0 & 0.3      & 500  & 200  & 42 \\
$\SM$    & C/N-scalar / EW    & 4     & 6.0 / 0.0 & 0.1--0.6 & 500  & 200  & 77 \\
$\SM$    & FSS (P1), $L=4$       & 4     & 6.0       & 0.2      & 500  & 500  & 99 \\
$\SM$    & FSS (P1), $L=6$       & 6     & 6.0       & 0.2      & 800  & 500  & 100 \\
$\SM$    & FSS (P1), $L=8$       & 8     & 6.0       & 0.2      & 1000 & 500  & 101 \\
$\SM$    & FSS (P1), $L=10$      & 10    & 6.0       & 0.2      & 1200 & 500  & 102 \\
\bottomrule
\end{tabular}
\end{table}

%%==============================================================
\section{E7 Gene-Network Positive-Control Protocol}
\label{app:geneprotocol}
%%==============================================================

This appendix documents the protocol used to produce the synthetic positive-control value
$\bar{\R}=0.345\pm0.024$ (s.d.) over 30 Alon-type synthetic networks.
The definitive empirical measurement on the real \textit{E.~coli} K-12
TF network is reported in Appendix~\ref{app:regulondb}.

\subsection*{Network Construction}

We generated synthetic regulatory networks following the structural conventions
in \citet{Alon2007}: directed graphs with $N\in\{16,32,64\}$ nodes, signed edges
(activation $=+1$, repression $=-1$), mean in-degree $\langle k\rangle\approx2$,
and no self-loops.  Ten random instances were drawn at each size ($N_\text{nets}=30$
total).  Edge signs were drawn Bernoulli($p=0.35$) for activation, matching the
$\sim$35\% activation fraction reported for the \textit{E.~coli}
transcription-factor network.

\subsection*{BFS Path Assignment and MIXED Fraction}

The standard E7 BFS protocol (§\ref{sec:framework}) is applied verbatim.
BFS from root node $r=0$ defines a spanning tree.  For each triangle
$(i,j,k)$ with all three directed edges present, the sign product is
\begin{equation}
  s_{ijk} = \mathrm{sgn}(e_{ij}\cdot e_{jk}\cdot e_{ki})\in\{+1,-1\},
\end{equation}
where $e_{ij}\in\{+1,-1\}$ is the signed adjacency weight.  The triangle is MIXED
if $s_{ijk}=-1$.  The MIXED fraction
$\R = \#\text{MIXED}/\#\text{triangles}$ uses the same formula as §\ref{sec:observables}.

\subsection*{Observed Values and Caveats}

Across the 30 synthetic networks we obtained $\R\in[0.32,0.48]$ with mean
$\bar{\R}=0.345\pm0.024$ (s.d.).  This is the source of the ``biological range''
cited in §6--8.

\textbf{Two important caveats:}
\begin{enumerate}
  \item The networks are \emph{synthetic}, not real biological data.  The range
    $[0.32,0.48]$ reflects the Alon-type graph topology, not a measurement
    on RegulonDB or YEASTRACT.
  \item The E7 BFS protocol was used for \emph{both} gauge-theory and gene-network
    computations, so the comparison is method-consistent but not independent.
    The protocol has now been applied to the real \textit{E.~coli} K-12 network
    from RegulonDB~\citep{SantosZavaleta2019}; results are reported in Appendix~\ref{app:regulondb}.
\end{enumerate}

%%==============================================================
\section{Open Items for Future Revision}
\label{app:openitems}
%%==============================================================

\begin{itemize}
  \item \textbf{EW-only finite-size check (Task F):}
    Run EW-only ($\beta_3=0$) scans at $L=6,8$ to test whether the gap
    EW-only $\R_q >$ full-SM $\R_q$ (Table~\ref{tab:sm_qvsl}) persists
    under FSS or closes.  A closing gap confirms the finite-$L$ mechanism
    sketched in §\ref{sec:sm_qvsl}; a persistent gap requires a revised
    theoretical explanation.
  \item \textbf{YEASTRACT empirical comparison:}
    Compute $\R$ on the signed \textit{S.~cerevisiae} TF network from
    YEASTRACT~\citep{Monteiro2020} using the E7 BFS protocol;
    compare with E.~coli RegulonDB result (Appendix~\ref{app:regulondb}).
  \item \textbf{Analytical commutator norm (Task N):}
    Derive $\mathcal{C}(\SU{2})$ and $\mathcal{C}(\SU{3})$ analytically
    from the Killing form or structure constants.
\end{itemize}

%%==============================================================
\section{RegulonDB MIXED-Fraction Computation}
\label{app:regulondb}
%%==============================================================

\subsection*{Data Source}

Network data were downloaded from RegulonDB v11.2~\citep{SantosZavaleta2019}
(\textit{E.~coli} K-12):
\texttt{NetWorkTFGene.txt} (6{,}061 TF\,$\to$\,gene interactions) and
\texttt{NetWorkTFTF.txt} (583 TF\,$\to$\,TF interactions).
Both files provide the direction of regulation ($+$ activation, $-$ repression,
$?$ unknown) and a confidence level (Confirmed, Strong, Weak).

\subsection*{Graph Construction}

Interactions with unknown sign ($?$) were excluded.
For each unordered node pair we collected all reported signs; if both $+$ and $-$
were present for the same pair (conflict), that edge was excluded (conservative).
Self-loops were dropped.
The E7 BFS MIXED protocol (Appendix~\ref{app:geneprotocol}) was applied to the
resulting unsigned undirected signed graph.
\textbf{Root independence:} on an undirected signed graph, the triangle sign
product $s_{ijk}=e_{ij}\cdot e_{jk}\cdot e_{ki}$ uses all three edge weights
directly, independently of spanning-tree structure; $\R$ is therefore root-independent.
This was confirmed by repeating the computation with 10 random root choices;
variation was $|\Delta\R|<0.001$ in both subsets.
Two confidence subsets were analysed:

\begin{itemize}
  \item \textbf{All signed} (Confirmed + Strong + Weak): 2{,}559 nodes,
    5{,}397 edges, 273 conflict-dropped pairs.
  \item \textbf{High-confidence} (Confirmed + Strong only): 2{,}137 nodes,
    3{,}666 edges, 183 conflict-dropped pairs.
\end{itemize}

\subsection*{Null Distribution}

For each subset, $N_\text{boot}=2{,}000$ null replicates were generated by
uniformly permuting all edge signs (preserving the total edge set), and
$\R_\text{null}$ was re-computed on each replicate.

\subsection*{Results}

\begin{table}[H]
\centering
\caption{RegulonDB \textit{E.~coli} K-12 MIXED-fraction results.
  $\R_\text{obs}$ is the directly computed value (bootstrap CI from triangle
  resampling). $\R_\text{null}$ is the sign-permutation null.}
\label{tab:regulondb}
\small
\begin{tabular}{lrrrrr}
\toprule
Subset & $N_\text{tri}$ & $\R_\text{obs}$ & $\R_\text{null}$ & $z$ & $p$ (2-tailed) \\
\midrule
All signed       & 1{,}559 & $0.413\pm0.012$ & $0.501\pm0.013$ & $-6.98\sigma$ & $<0.001$ \\
High-conf.       &   725   & $0.349\pm0.018$ & $0.500\pm0.019$ & $-8.07\sigma$ & $<0.001$ \\
\bottomrule
\end{tabular}
\end{table}

Both subsets give $\R_\text{obs}$ significantly \emph{below} the random null
($z\approx{-7}$ to $-8\sigma$, $p<0.001$, $N_\text{boot}=2{,}000$).
This confirms that the \textit{E.~coli} TF network is enriched for
\emph{balanced} (non-MIXED) triangles, consistent with network balance
theory~\citep{Alon2007}.

The high-confidence value $\R_\text{RegulonDB}=0.349\pm0.018$ overlaps
numerically with the SU(3) FSS extrapolation
$R_\infty^{\SU{3}}=0.3539\pm0.0195$ from \S\ref{sec:smfss}.
We do \emph{not} interpret this overlap as a causal or mechanistic connection;
neither the gauge algebra nor the gene network imposes a common fixed point.
The overlap is noted as a numerical coincidence warranting further
investigation (see \S\ref{sec:conclusions}).

\subsection*{Reproducibility}

Analysis script: \texttt{daft-forward/E1E4E7/h\_regulondb.py}.
Random seed: 777.
Output: \texttt{h\_regulondb\_results.json} (deposited alongside simulation data).

%%==============================================================
\section{Commutator Norm $\mathcal{C}(G)$: Monte Carlo Computation}
\label{app:commutator}
%%==============================================================

\subsection*{Definition and Method}

The commutator norm $\mathcal{C}(G) = \mathbb{E}_{U_1,U_2\sim\mu_G}[\|[U_1,U_2]\|_F]$
(§\ref{sec:framework}) was evaluated by Monte Carlo integration.
Haar-random matrices were generated via the QR method: draw
$X\sim\mathcal{CN}(0,I)$, compute $X=QR$, normalise column phases of $Q$ via
$Q \leftarrow Q\,\mathrm{diag}(R_{ii}/|R_{ii}|)$, then enforce $\det Q = 1$.
Both $U_1$ and $U_2$ were drawn independently, and
$\|[U_1,U_2]\|_F = \|U_1 U_2 - U_2 U_1\|_F$ was computed for each pair.

\subsection*{Results}

\begin{table}[H]
\centering
\caption{Commutator norm $\mathcal{C}(G)$ by Monte Carlo ($N_\text{pairs}=200{,}000$,
  seed 314). Values for the SM product group are not computed (multi-sector transport).}
\label{tab:commutator}
\begin{tabular}{llll}
\toprule
Group & $\mathcal{C}(G)$ & 95\% CI & Method \\
\midrule
$\UN{1}$ & $0$ (exact) & --- & Abelian \\
$\SU{2}$ & $1.601\pm0.002$ & $[1.598, 1.604]$ & Monte Carlo \\
$\SU{3}$ & $2.275\pm0.001$ & $[2.274, 2.277]$ & Monte Carlo \\
\bottomrule
\end{tabular}
\end{table}

\subsection*{Note on Previously Stated Values}

The original manuscript stated $\mathcal{C}(\SU{2})\approx2.3$ and
$\mathcal{C}(\SU{3})\approx2.633$ (sourced from prior informal estimates).
The recomputed values are $1.601$ and $2.275$ respectively.
The \emph{ordering} $\mathcal{C}(\UN{1})<\mathcal{C}(\SU{2})<\mathcal{C}(\SU{3})$
is confirmed.
The corrected ratio is $\mathcal{C}(\SU{3})/\mathcal{C}(\SU{2})=1.421$, compared
to the simulated $\R$ ratio $R_\infty^{\SU{3}}/R_\infty^{\SU{2}}=1.176\pm0.015$;
this indicates a saturating (sub-linear) dependence of $\R$ on $\mathcal{C}(G)$
rather than strict proportionality.

\subsection*{Reproducibility}

Script: \texttt{daft-forward/E1E4E7/n\_commutator.py}.
Output: \texttt{n\_commutator\_results.json}.

%%==============================================================
\section{ER Random Graph Null Distribution for $\R$}
\label{app:null_er}
%%==============================================================

\subsection*{Protocol}

To assess whether the range $\R\in[0.32,0.48]$ is easily achieved by structurally
random signed graphs, we generated $N_\text{nets}=10{,}000$ Erd{\H o}s--R{\'e}nyi
random signed graphs per network size $N\in\{16,32,64\}$, matching the
positive-control gene networks (Appendix~\ref{app:geneprotocol}) in:
\begin{itemize}
  \item Edge probability: $p_\text{edge}=2/(N-1)$ (mean undirected degree $\approx2$).
  \item Sign: Bernoulli$(p=0.35)$ activation ($+1$), else repression ($-1$).
\end{itemize}
The E7 MIXED fraction was computed on each graph (graphs without any triangle
were excluded from the distribution: $\sim$30--35\% of networks at these densities).

\subsection*{Results}

\begin{table}[H]
\centering
\caption{ER random signed graph null distribution for $\R$.
  ``no-tri fraction'' = fraction of networks with zero triangles (excluded from $\R$).}
\label{tab:er_null}
\small
\begin{tabular}{lrrrrr}
\toprule
$N$ & $\bar{\R}$ & $\sigma_\R$ & 95\%~CI & $P(\R\in[0.32,0.48])$ & no-tri \\
\midrule
16  & $0.517$ & $0.408$ & $[0.00,\,1.00]$ & $7.1\%$ & $34.9\%$ \\
32  & $0.516$ & $0.412$ & $[0.00,\,1.00]$ & $5.9\%$ & $31.5\%$ \\
64  & $0.516$ & $0.417$ & $[0.00,\,1.00]$ & $5.9\%$ & $28.7\%$ \\
Pooled & $0.516$ & $0.412$ & $--$ & $6.3\%$ & $--$ \\
\bottomrule
\end{tabular}
\end{table}

\subsection*{Interpretation}

The ER null distribution is approximately bimodal: either no triangles exist (excluded
as NaN) or $\R\approx0.5$ (random sign product).
Only $6.3\%$ of random ER graphs land in $[0.32,0.48]$, confirming:
\begin{enumerate}
  \item The range is not trivially populated by sparse random graphs.
  \item The gene-network controls ($\R\approx0.35$, Appendix~\ref{app:geneprotocol})
        achieve this range because of their structured (non-random) topology, not by
        chance.
  \item The SU(3) MIXED fraction $R_\infty^{\SU{3}}=0.354\pm0.020$, which falls in
        this range, is not what one would expect from a structurally random gauge
        network.
\end{enumerate}

\subsection*{Reproducibility}

Script: \texttt{daft-forward/E1E4E7/p\_null\_distribution.py}.
Output: \texttt{p\_null\_distribution\_results.json}.

\end{document}
